\textbf{章节结构}：
\begin{itemize}
    \item \S 6.1 稳定性
    \begin{itemize}
        \item 6.1.1 李雅普诺夫稳定性
        \item 6.1.2 按线性近似决定稳定性
        \item 6.1.3 V 函数方法（李雅普诺夫第二方法）
    \end{itemize}
    \item \S 6.2 定性
    \begin{itemize}
        \item 6.2.1 奇点
        \item 6.2.2 极限环
        \item 6.2.3 平面图貌
    \end{itemize}
    \item \S 6.3 哈密顿系统、孤立子和混沌
\end{itemize}

\section{稳定性}

\subsection{引入例：一阶非线性方程的解的图貌}

求 $\dfrac{dy}{dt} = Ay - By^2$ 的解的图貌，其中 $A, B$ 为常数且 $AB > 0$，初值 $y(0) = y_0$。

解：$y(t) = \dfrac{A}{B + Ce^{-At}}$，平衡解 $y_1(t) = 0, \ y_2(t) = \dfrac{A}{B}$。满足初值的解：
$$y(t) = \frac{A}{B + \left(\frac{A}{y_0} - B\right)e^{-At}}$$

\begin{itemize}
    \item \textbf{(a)} $A > 0, B > 0$：$y_0 > 0 \Rightarrow y(t) \to \dfrac{A}{B}$（$y_2$ 稳定）；$y_0 < 0 \Rightarrow y(t) \to -\infty$（发散）
    \item \textbf{(b)} $A < 0, B < 0$：$y_0 > \dfrac{A}{B} \Rightarrow y(t) \to \infty$；$y_0 < \dfrac{A}{B} \Rightarrow y(t) \to 0$（$y_1$ 稳定）
\end{itemize}

\textbf{约定}：当 $A > 0, B > 0$ 时 $y_2(t)$ 稳定、$y_1(t)$ 不稳定；当 $A < 0, B < 0$ 时 $y_1(t)$ 稳定、$y_2(t)$ 不稳定。

\textbf{化为零解稳定性的平移}：一般非线性方程组 $\dot{\mathbf{y}} = \mathbf{g}(t, \mathbf{y})$ 的特解 $\mathbf{y} = \boldsymbol{\varphi}(t)$，用平移变换 $\mathbf{x} = \mathbf{y} - \boldsymbol{\varphi}(t)$ 化为
$$\dot{\mathbf{x}} = \mathbf{f}(t, \mathbf{x}), \quad \mathbf{f}(t, 0) = 0$$
其中 $\mathbf{f}(t, \mathbf{x}) = \mathbf{g}(t, \boldsymbol{\varphi}(t) + \mathbf{x}) - \dot{\boldsymbol{\varphi}}(t)$。原特解稳定 $\iff$ 变换后零解 $\mathbf{x} \equiv 0$ 稳定。

\textbf{驻定方程与平衡解}：右端不含自变量的微分方程称为\textbf{驻定（自治）微分方程}；右端为零得到的代数方程的解即常数特解，称为\textbf{驻定解}或\textbf{平衡解}（如 $Ay - By^2 = 0$ 的解 $y_1 = 0, y_2 = A/B$）。

\subsection{李雅普诺夫稳定性}

考虑 $\dfrac{d\mathbf{x}}{dt} = \mathbf{f}(t, \mathbf{x}), \ \mathbf{f}(t, 0) = 0$ (5)，假设 $\mathbf{f}$ 在包含原点的域 $G$ 内具有连续偏导数（满足存在唯一性、延拓性、连续可微性）。

\begin{itemize}
    \item \textbf{稳定}：若对任意 $\varepsilon > 0$，存在 $\delta = \delta(\varepsilon, t_0) > 0$，使对任意初值 $\mathbf{x}_0$，只要 $\|\mathbf{x}_0\| \le \delta$，则解 $\mathbf{x}(t)$ 在一切 $t \ge t_0$ 上都有 $\|\mathbf{x}(t)\| \le \varepsilon$，称零解 $\mathbf{x} \equiv 0$ \textbf{稳定}（Lyapunov 稳定）
    \item \textbf{渐近稳定}：零解稳定，且存在 $\delta_0 > 0$，使当 $\|\mathbf{x}_0\| \le \delta_0$ 时对应解满足 $\lim_{t\to\infty}\mathbf{x}(t) = 0$
    \item \textbf{稳定域（吸引域）}：若零解渐近稳定，且存在区域 $D_0$，当且仅当 $\mathbf{x}_0 \in D_0$ 时解满足 $\lim_{t\to\infty}\mathbf{x}(t) = 0$
    \item \textbf{全局渐近稳定}：稳定域为整个空间（$\delta_0 = +\infty$）
    \item \textbf{不稳定}：存在 $\varepsilon > 0$，无论 $\delta$ 多小，总可找到 $\|\mathbf{x}_0\| \le \delta$ 使解在某个 $t_1 > t_0$ 满足 $\|\mathbf{x}(t_1)\| = \varepsilon$
\end{itemize}

\textbf{平面上三种稳定性态}：
\begin{itemize}
    \item[(a)] \textbf{稳定}：轨线始终保持在 $\varepsilon$ 邻域内
    \item[(b)] \textbf{渐近稳定}：轨线不仅不离开小邻域，且最终收敛到原点
    \item[(c)] \textbf{不稳定}：任意小扰动都会使轨线最终离开 $\varepsilon$ 邻域
\end{itemize}

\textbf{例（logistic 型）}：$\dfrac{dy}{dt} = Ay - By^2$
\begin{itemize}
    \item 情形一 $A < 0, B < 0$：零解 $y=0$ 渐近稳定，稳定域为 $y < A/B$，$y_2 = A/B$ 不稳定
    \item 情形二 $A > 0, B > 0$：$y_2 = A/B$ 渐近稳定，稳定域为 $y > 0$，零解不稳定
\end{itemize}

\subsubsection{常系数线性微分方程组的稳定性}

$\dfrac{d\mathbf{x}}{dt} = \mathbf{A}\mathbf{x}$ (8)，特征方程 $\det(\mathbf{A} - \lambda\mathbf{E}) = 0$ (10)，一般解 $x(t) = \sum_i \sum_{m=0}^{l_i} c_{im}t^m e^{\lambda_i t}$。

\begin{myTheorem}{稳定性判据}
\begin{itemize}
    \item 若全部特征根 $\mathrm{Re}(\lambda_i) < 0$，零解\textbf{渐近稳定}
    \item 若存在特征根 $\mathrm{Re}(\lambda_i) > 0$，零解\textbf{不稳定}
    \item 若无正实部根，但存在零根或 $\mathrm{Re}(\lambda_i) = 0$ 的根，零解\textbf{可能稳定也可能不稳定}，需根据零根的重数与结构进一步判断
\end{itemize}
\end{myTheorem}

\subsection{按线性近似决定稳定性}

考虑 $\dfrac{d\mathbf{x}}{dt} = \mathbf{A}\mathbf{x} + \mathbf{R}(\mathbf{x}), \ \mathbf{R}(0) = 0$ (11)，假设 $\lim_{\|\mathbf{x}\|\to 0}\dfrac{\|\mathbf{R}(\mathbf{x})\|}{\|\mathbf{x}\|} = 0$（非线性项为高阶小量，即微小非线性扰动）。

\begin{myTheorem}{线性化稳定性定理}
\begin{itemize}
    \item 若特征方程 (10) 无零根或零实部特征根，则非线性系统 (11) 与线性系统 (8) 的零解稳定性完全一致：
    \begin{itemize}
        \item 所有特征根实部 $< 0$ $\Rightarrow$ 零解渐近稳定
        \item 存在特征根实部 $> 0$ $\Rightarrow$ 零解不稳定
    \end{itemize}
    \item 若存在零根或零实部特征根，则为\textbf{临界情形}：线性系统 (8) 不能决定 (11) 零解的稳定性，必须进一步分析高阶项 $\mathbf{R}(\mathbf{x})$ 的作用
\end{itemize}
\end{myTheorem}

\begin{example}[有阻力的数学摆]
\begin{equation}\frac{d^2\varphi}{dt^2} + \frac{\mu}{m}\frac{d\varphi}{dt} + \frac{g}{l}\sin\varphi = 0 \tag{6.15}\end{equation}
令 $x = \varphi, y = \dfrac{d\varphi}{dt}$，化为一阶方程组：
\begin{equation}\frac{dx}{dt} = y, \qquad \frac{dy}{dt} = -\frac{g}{l}\sin x - \frac{\mu}{m}y \tag{6.16}\end{equation}
\end{example}

\begin{itemize}
    \item 原点 $(0,0)$ 为零解。分解线性 $+$ 非线性：$\dfrac{dy}{dt} = -\dfrac{g}{l}x - \dfrac{\mu}{m}y - \dfrac{g}{l}(\sin x - x)$，非线性项 $R(x,y) = -\dfrac{g}{l}(\sin x - x) = \dfrac{g}{l}\left(\dfrac{x^3}{3!} - \dfrac{x^5}{5!} + \cdots\right) = o(\|(x,y)\|)$
    \item 线性近似 (6.17) 的特征方程 $\lambda^2 + \dfrac{\mu}{m}\lambda + \dfrac{g}{l} = 0$，$\lambda_{1,2} = -\dfrac{\mu}{2m} \pm \dfrac{1}{2}\sqrt{\left(\dfrac{\mu}{m}\right)^2 - \dfrac{4g}{l}}$
    \item 因 $\mu > 0$，$\mathrm{Re}(\lambda_{1,2}) < 0$，由定理 2 原点零解\textbf{渐近稳定}
\end{itemize}

\textbf{其他驻定解}：$y = 0, \sin x = 0 \Rightarrow x = n\pi$。对代表性驻定解 $(\pi, 0)$ 平移 $x^* = x - \pi, y^* = y$，线性化得 $\dfrac{dx^*}{dt} = y^*, \ \dfrac{dy^*}{dt} = \dfrac{g}{l}x^* - \dfrac{\mu}{m}y^*$，特征根 $\lambda_{1,2} = -\dfrac{\mu}{2m} \pm \dfrac{1}{2}\sqrt{\left(\dfrac{\mu}{m}\right)^2 + \dfrac{4g}{l}}$，一正一负，为\textbf{鞍点}，\textbf{不稳定}。

\textbf{无阻力情形（$\mu = 0$）}：线性系统特征根 $\lambda_{1,2} = \pm i\sqrt{g/l}$ 为纯虚根，线性系统零解稳定；非线性系统在原点的稳定性属\textbf{临界情形}，不能仅靠线性化判定。但平衡解 $x = \pi, y = 0$ 依然不稳定。

\subsubsection{霍维茨（Hurwitz）判别法}

$n$ 次方程 $a_0\lambda^n + a_1\lambda^{n-1} + \cdots + a_{n-1}\lambda + a_n = 0$（$a_0 > 0$），定义 Hurwitz 行列式：
$$\Delta_1 = a_1, \quad \Delta_2 = \begin{vmatrix} a_1 & a_0 \\ a_3 & a_2 \end{vmatrix}, \quad \Delta_3 = \begin{vmatrix} a_1 & a_0 & 0 \\ a_3 & a_2 & a_1 \\ a_5 & a_4 & a_3 \end{vmatrix}, \ldots, \Delta_n = a_n\Delta_{n-1}$$
其中 $a_i = 0$（$i > n$）。

\begin{myTheorem}{Hurwitz 判据}
方程全部根具有负实部的充要条件：$a_1 > 0, \Delta_2 > 0, \Delta_3 > 0, \ldots, \Delta_{n-1} > 0, a_n > 0$。
\end{myTheorem}

\begin{example}
三维系统
$$\begin{cases} \dot{x} = -2x + y - z + x^2e^x \\ \dot{y} = x - y + x^3y + z^2 \\ \dot{z} = x + y - z - e^x(y^2 + z^2) \end{cases}$$
线性化 Jacobian $\mathbf{A} = \begin{pmatrix} -2 & 1 & -1 \\ 1 & -1 & 0 \\ 1 & 1 & -1 \end{pmatrix}$，特征方程 $\lambda^3 + 4\lambda^2 + 5\lambda + 3 = 0$。$a_1 = 4 > 0$，$\Delta_2 = \begin{vmatrix} 4 & 1 \\ 3 & 5 \end{vmatrix} = 17 > 0$，$a_3 = 3 > 0$。由 Hurwitz 判据全部特征根实部为负，故原系统零解\textbf{渐近稳定}。
\end{example}

\begin{example}
三次方程 $\lambda^3 + (a+b+1)\lambda^2 + b(a+c)\lambda + 2ab(c-1) = 0$（$a, b, c > 0$）根均具负实部的条件：
$$a_1 = a+b+1 > 0, \quad \Delta_2 = (a+b+1)(b(a+c) - 2ab(c-1)) > 0, \quad a_3 = 2ab(c-1) > 0 \iff c > 1$$
由 $\Delta_2 > 0$ 得参数区域：$a < b+1, c > 1$；或 $a > b+1, 1 < c < c_0$，其中 $c_0 = \dfrac{a(a+b+3)}{a-b-1}$。
\end{example}

\subsection{V 函数方法（李雅普诺夫第二方法）}

考虑 $\dfrac{d\mathbf{x}}{dt} = \mathbf{f}(\mathbf{x}), \ \mathbf{f}(0) = 0, \ \mathbf{x} \in \mathbb{R}^n$ (14)，$\mathbf{f}$ 在 $G = \{\|\mathbf{x}\| \le A\}$ 内有连续偏导数。

\textbf{V 函数的定义}：$V(\mathbf{x})$ 在 $\|\mathbf{x}\| \le H$ 内连续可微，$V(0) = 0$。
\begin{itemize}
    \item $V(\mathbf{x}) \ge 0$ 称为\textbf{常正}
    \item $\mathbf{x} \neq 0$ 时 $V(\mathbf{x}) > 0$ 称为\textbf{定正}
    \item 若 $-V$ 为定正（常正），则 $V$ 为\textbf{定负（常负）}
\end{itemize}

\textbf{沿解的全导数}：
$$\frac{dV}{dt} = \sum_{i=1}^{n}\frac{\partial V}{\partial x_i}\frac{dx_i}{dt} = \sum_{i=1}^{n}\frac{\partial V}{\partial x_i}f_i(\mathbf{x}) = \frac{\partial V}{\partial \mathbf{x}}\cdot\mathbf{f}(\mathbf{x})$$

\begin{example}
$V(x,y) = (x+y)^2$ 常正；$V = (x+y)^2 + y^2$ 定正；$V = \sin(x+y)^2$ 在 $x^2+y^2 < \pi$ 上定正、全平面变号；二次型 $V = ax^2 + bxy + cy^2$ 当 $a > 0, 4ac - b^2 > 0$ 时定正。
\end{example}

\begin{myTheorem}{Lyapunov 稳定与渐近稳定}
\begin{itemize}
    \item 若存在在原点邻域\textbf{定正}的函数 $V(\mathbf{x})$，且沿系统 (14) 的全导数 $\dfrac{dV}{dt}$ 为\textbf{常负函数}或恒为零，则零解\textbf{稳定}
    \item 若存在定正函数 $V(\mathbf{x})$，其全导数为\textbf{定负函数}，则零解\textbf{渐近稳定}
\end{itemize}
\end{myTheorem}

\textbf{不稳定判据}：若存在函数 $V(\mathbf{x})$ 及常数 $\mu \ge 0$，使得 $\dfrac{dV}{dt} = \mu V(\mathbf{x}) + W(\mathbf{x})$，且满足：
\begin{itemize}
    \item $\mu = 0$ 时，$W(\mathbf{x})$ 为定正函数
    \item $\mu \neq 0$ 时，$W(\mathbf{x})$ 为常负函数或恒为零
    \item 在 $\mathbf{x} = 0$ 的任意小邻域内至少存在一点 $\mathbf{x}$ 使 $V(\mathbf{x}) > 0$
\end{itemize}
则系统 (14) 的零解\textbf{不稳定}。

\textbf{证明思路}：
\begin{itemize}
    \item \textbf{稳定性}：任取 $\varepsilon < H$，因 $V$ 定正，环带 $\varepsilon \le \|\mathbf{x}\| \le H$ 上 $V$ 有正下界 $l$；又 $V(0) = 0$ 连续，存在 $\delta$ 使 $\|\mathbf{x}_0\| \le \delta \Rightarrow V(\mathbf{x}_0) < l$。沿解 $\dfrac{dV}{dt} \le 0$，故 $V(\mathbf{x}(t)) \le V(\mathbf{x}_0) < l$，反设 $\|\mathbf{x}(t^*)\| = \varepsilon$ 则 $V(\mathbf{x}(t^*)) \ge l$ 矛盾
    \item \textbf{渐近稳定性}：$V$ 沿解严格递减且有下界，$\lim V(\mathbf{x}(t)) = c \ge 0$。若 $c > 0$ 则轨线在环带内，$\dfrac{dV}{dt} \le m < 0$，积分得 $V \to -\infty$ 矛盾；故 $V \to 0$，进而 $\mathbf{x}(t) \to 0$
    \item \textbf{不稳定性}：由 $\dfrac{dV}{dt} \ge \mu V$，乘 $e^{-\mu(t-t_0)}$ 积分得 $V(\mathbf{x}(t)) \ge V(\mathbf{x}_0)e^{\mu(t-t_0)}$，$\mu > 0$ 时 $V \to \infty$ 与 $\|\mathbf{x}\| \le H$ 下有界矛盾
\end{itemize}

\textbf{几何解释}：等值面 $V(\mathbf{x}) = c$ 是围绕原点的闭曲面。零解稳定 $\Rightarrow$ 轨线停留在某个闭曲面内；渐近稳定 $\Rightarrow$ 轨线沿一族闭曲面逐层向原点收缩。

\begin{example}
$\dot{x} = -y + ax^3, \ \dot{y} = x + ay^3$。线性化 $\dot{x} = -y, \dot{y} = x$，特征根 $\lambda = \pm i$ 属\textbf{临界}，线性化无法判定。取 $V = \dfrac{1}{2}(x^2 + y^2)$，沿解 $\dot{V} = x\dot{x} + y\dot{y} = a(x^4 + y^4)$：
\begin{itemize}
    \item $a < 0$：$\dot{V}$ 定负，零解\textbf{渐近稳定}
    \item $a > 0$：$\dot{V} > 0$ 定正，零解\textbf{不稳定}
    \item $a = 0$：$\dot{V} \equiv 0$，$V$ 保持常数，零解\textbf{稳定但非渐近稳定}
\end{itemize}
\end{example}

\begin{myTheorem}{Lyapunov 推广}
若存在定正 $V(\mathbf{x})$，其全导数 $\dot{V}$ 为常负，但使 $\dot{V}(\mathbf{x}) = 0$ 的点集除 $\mathbf{x} = 0$ 外不包含整条正半轨线，则零解\textbf{渐近稳定}。
\end{myTheorem}

\begin{remark}
直观：虽然 $\dot{V} \le 0$ 允许某些点 $\dot{V} = 0$，但因这些点不含任何完整正半轨线，轨线不能"停在"某个闭曲线 $V = c$ 上。
\end{remark}

\begin{example}[含阻力数学摆（V 函数法）]
$V(x,y) = \dfrac{1}{2}y^2 + \dfrac{g}{l}(1 - \cos x)$ 原点附近定正，沿解 $\dfrac{dV}{dt} = -\dfrac{\mu}{m}y^2 \le 0$。
\begin{itemize}
    \item $\mu = 0$：$\dot{V} = 0$，$V$ 为第一积分，零解\textbf{稳定}（临界情形）
    \item $\mu > 0$：$\dot{V} = -\dfrac{\mu}{m}y^2$ 常负，但 $\dot{V} = 0$ 的集合是直线 $y = 0$，除 $(0,0)$ 外不包含整条正半轨线，由定理 5 零解\textbf{渐近稳定}
\end{itemize}
\end{example}

\textbf{李雅普诺夫稳定性理论概述}：第一方法（线性化/特征值判别）与第二方法（V 函数）。第二方法不依赖解的显式表达，是非线性系统研究最重要的工具。核心难点是\textbf{临界情形}（零实部特征值），需依赖非线性高次项判定。后续推广：全局稳定性（无限大 Lyapunov 函数）、时滞/泛函微分方程、脉冲系统、时标动态方程。

\textbf{作业（6.1）}：p157-158：1（2）、3（2、3、5）、4（1、2、3）、5（1、2、3）、6

\section{定性理论}

考虑平面一阶方程组 $\begin{cases} \dot{x} = X(x,y) \\ \dot{y} = Y(x,y) \end{cases}$ (6.15)，$X, Y$ 对 $(x,y)$ 有连续偏导数。解 $(x(t), y(t))$ 在三维空间 $Otxy$ 中形成\textbf{积分曲线}；平面 $Oxy$ 称为\textbf{相平面}，积分曲线在相平面上的投影称为\textbf{轨线（轨道）}。

若 $X^2 + Y^2$ 不恒为零，可化为一阶方程 $\dfrac{dy}{dx} = \dfrac{Y(x,y)}{X(x,y)}$（$X \neq 0$）或 $\dfrac{dx}{dy} = \dfrac{X(x,y)}{Y(x,y)}$（$Y \neq 0$）。相平面上轨线\textbf{不能相交}（同一初始点只对应唯一一条轨线）。

\subsection{奇点}

若常数 $x = x^*, y = y^*$ 同时满足 $X(x^*, y^*) = 0, \ Y(x^*, y^*) = 0$，则 $(x^*, y^*)$ 给出方程组的一个\textbf{驻定解}，相平面上对应点称为\textbf{奇点}（平衡点、平衡态）。通过平移可将奇点移到原点。

在奇点邻域取线性项，得线性近似方程组
\begin{equation}\frac{dx}{dt} = ax + by, \qquad \frac{dy}{dt} = cx + dy \tag{6.18}\end{equation}
设 $\det\mathbf{A} = \begin{vmatrix} a & b \\ c & d \end{vmatrix} \neq 0$，则原点是 (6.18) 的唯一奇点。特征方程
$$\begin{vmatrix} a - \lambda & b \\ c & d - \lambda \end{vmatrix} = 0 \iff \lambda^2 + p\lambda + q = 0, \quad p = -(a+d), \ q = ad - bc, \ \Delta = p^2 - 4q$$
通过非奇异实线性变换可将系数矩阵化为四种标准形之一：$\begin{pmatrix} \lambda & 0 \\ 0 & \mu \end{pmatrix}$、$\begin{pmatrix} \lambda & 1 \\ 0 & \lambda \end{pmatrix}$、$\begin{pmatrix} \lambda & 0 \\ 0 & \lambda \end{pmatrix}$、$\begin{pmatrix} \alpha & \beta \\ -\beta & \alpha \end{pmatrix}$，对应不同类型奇点。

\subsubsection{五种情形的奇点分类}

\begin{itemize}
    \item \textbf{情形 I（结点，$q > 0, \Delta > 0$）}：同号相异实根 $\lambda_1 \neq \lambda_2$。标准形式 $\dot{\xi} = \lambda_1\xi, \dot{\eta} = \lambda_2\eta$，通解 $\xi = Ae^{\lambda_1t}, \eta = Be^{\lambda_2t}$。若 $\lambda_1 > \lambda_2$，轨线斜率 $\dfrac{\eta}{\xi} = \dfrac{B}{A}e^{(\lambda_2-\lambda_1)t} \to 0$，轨线在原点\textbf{切 $\xi$ 轴}。特征方向上的轴为轨线。\textbf{稳定结点}（$\lambda_{1,2} < 0$，渐近稳定）；\textbf{不稳定结点}（$\lambda_{1,2} > 0$）
    \item \textbf{情形 II（鞍点，$q < 0$）}：异号实根 $\lambda_1\lambda_2 < 0$。仅沿两条特征方向的四条轨线趋于/离开原点，其他轨线"穿过"原点附近，为典型\textbf{不稳定}平衡点
    \item \textbf{情形 III（重根，$\Delta = 0$）}：
    \begin{itemize}
        \item 一般情形（$b \neq 0$ 或 $c \neq 0$）：标准形式 $\dot{\xi} = \lambda\xi + \eta, \dot{\eta} = \lambda\eta$，通解 $\xi = (At+B)e^{\lambda t}, \eta = Ae^{\lambda t}$，轨线切 $\xi$ 轴，为\textbf{退化结点}（$\lambda < 0$ 稳定，$\lambda > 0$ 不稳定）
        \item $b = c = 0$（$\lambda = a = d$）：通解 $x = Ae^{\lambda t}, y = Be^{\lambda t}$，所有轨线均为通过原点的半射线，为\textbf{奇结点}（$\lambda < 0$ 稳定，$\lambda > 0$ 不稳定）
    \end{itemize}
    \item \textbf{情形 IV（焦点，$\Delta < 0, p \neq 0$）}：共轭复根 $\lambda = \alpha \pm i\beta$。极坐标下 $\dot{r} = \alpha r, \dot{\theta} = -\beta$，通解 $r = Ae^{\alpha t}, \theta = -\beta t + B$，轨线为\textbf{对数螺旋线}。$\alpha < 0$ 稳定焦点（渐近稳定）；$\alpha > 0$ 不稳定焦点
    \item \textbf{情形 V（中心，$\Delta < 0, p = 0$）}：纯虚根 $\lambda = \pm i\beta$，轨线为同心圆，原点为\textbf{中心}，零解\textbf{稳定但非渐近稳定}
\end{itemize}

\begin{myTheorem}{线性奇点定理}
若 $\det\mathbf{A} \neq 0$，则奇点类型完全由特征方程 $\lambda^2 + p\lambda + q = 0$ 的根决定：
\begin{enumerate}
    \item 两个非零实根 $\lambda_1 \neq \lambda_2$：$\lambda_1\lambda_2 > 0$ 为结点（$\lambda < 0$ 稳定、$\lambda > 0$ 不稳定）；$\lambda_1\lambda_2 < 0$ 为鞍点（不稳定）
    \item 重根 $\lambda \neq 0$：一般情形退化结点；$b = c = 0$ 为奇结点；$\lambda < 0$ 稳定、$\lambda > 0$ 不稳定
    \item 共轭复根 $\lambda_{1,2} = \alpha \pm i\beta$：$\alpha \neq 0$ 为焦点（$\alpha < 0$ 稳定、$\alpha > 0$ 不稳定）；$\alpha = 0$ 为中心（稳定但非渐近稳定）
\end{enumerate}
\end{myTheorem}

\textbf{$(p,q)$ 平面划分}：分界曲线为抛物线 $p^2 - 4q = 0$，把结点、鞍点、焦点、中心等类型的区域分割开来。

\begin{example}
$\dfrac{d^2x}{dt^2} + 3\dfrac{dx}{dt} + 2x = 0$。令 $y = x'$ 得 $\begin{cases} \dot{x} = y \\ \dot{y} = -2x - 3y \end{cases}$，特征根 $\lambda_1 = -1, \lambda_2 = -2$ 为同号相异实根，原点为\textbf{稳定结点}，特征方向对应直线 $x + y = 0$ 与 $2x + y = 0$。
\end{example}

\subsection{极限环}

\begin{example}
$\begin{cases} \dot{x} = x + y - x(x^2+y^2) \\ \dot{y} = -x + y - y(x^2+y^2) \end{cases}$，极坐标下 $\dot{r} = r(1-r^2), \dot{\theta} = -1$：
\begin{itemize}
    \item $r \equiv 0$：原点奇点
    \item $r \equiv 1$：单位圆轨道（周期 $2\pi$ 的闭轨线，顺时针），是候选极限环
    \item $0 < R_1 < 1$：$\dot{r}|_{r=R_1} > 0$，轨线从内向外穿出；$R_2 > 1$：$\dot{r}|_{r=R_2} < 0$，轨线从外向内穿入
    \item 取环域 $D = \{R_1 < r < R_2\}$：轨线从两侧边界均流入 $D$，$D$ 内无奇点，故 $r = 1$ 为\textbf{吸引型（稳定）极限环}
\end{itemize}
\end{example}

\textbf{极限环的稳定性}：稳定（吸引型，邻域内轨线均向其靠拢）、不稳定（排斥型）、半稳定（一侧靠拢另一侧远离）。极限环必须是\textbf{孤立}闭轨线。

\begin{myTheorem}{环域定理}
设 $G$ 为方程组 (6.15) 的定义域，若 $G$ 内存在一个有界的环形闭域 $D$，满足：\textcircled{1} $D$ 内不含奇点；\textcircled{2} 经过 $D$ 上任一点的解轨线在 $t \ge t_0$（或 $t \le t_0$）时始终不离开 $D$。则必有二者之一成立：该解本身是周期解（闭轨线），或该解在正向（负向）时间上趋于 $D$ 内的某个周期解。
\end{myTheorem}

\begin{myTheorem}{Bendixson 判据}
若在单连通区域 $D^*$ 中 $\dfrac{\partial X}{\partial x} + \dfrac{\partial Y}{\partial y}$（散度）不变号，且在 $D^*$ 任何子域内不恒为零，则系统在 $D^*$ 内不存在周期解，也不存在极限环。
\end{myTheorem}

\begin{proof}
反证：假设存在周期轨 $\Gamma$ 围成区域 $D_\Gamma$，由 Green 公式 $\iint_{D_\Gamma}\mathrm{div}\,\mathbf{F}\,dx\,dy = \oint_\Gamma (X\,dy - Y\,dx) = \int_0^T (X\dot{y} - Y\dot{x})\,dt = 0$，但散度不变号且不恒为零时该积分 $\neq 0$，矛盾。
\end{proof}

\begin{example}
阻尼摆 $\begin{cases} \dot{x} = y \\ \dot{y} = -\dfrac{g}{l}\sin x - \dfrac{\mu}{m}y \end{cases}$，散度 $= -\dfrac{\mu}{m} < 0$ 不变号，由 Bendixson 判据\textbf{无周期解、无极限环}。
\end{example}

\textbf{Lienard 系统与 van der Pol 方程}：范德波尔方程 $x'' + \mu(x^2-1)x' + x = 0$ 属于更一般的 Lienard 系统 $x'' + f(x)x' + g(x) = 0$。定义 $F(x) = \int_0^x f(s)\,ds$，令 $y = x' + F(x)$，化为一阶系统 $\begin{cases} \dot{x} = y - F(x) \\ \dot{y} = -g(x) \end{cases}$。

\begin{myTheorem}{Lienard 方程定理}
假设
\begin{enumerate}
    \item $f(x), g(x)$ 对一切 $x$ 连续，$g(x)$ 满足局部 Lipschitz 条件
    \item $f(x)$ 为偶函数，$f(0) < 0$；$g(x)$ 为奇函数，当 $x \neq 0$ 时 $xg(x) > 0$
    \item 当 $x \to \pm\infty$ 时 $F(x) \to \pm\infty$，$F(x)$ 有唯一正零点 $x = a$，且当 $x \ge a$ 时 $F(x)$ 单调增加
\end{enumerate}
\textbf{结论}：Lienard 方程有唯一周期解，即方程组有一个（稳定的）极限环。
\end{myTheorem}

\textbf{van der Pol 方程}：$x'' + \mu(x^2-1)x' + x = 0$ 等价于 $\begin{cases} \dot{x} = y - \mu\left(\frac{x^3}{3} - 1\right) \\ \dot{y} = -x \end{cases}$，满足定理 9 条件，故存在一个极限环。随着 $\mu$ 由小变大，其极限环由接近圆形逐渐变为复杂形状。

\subsection{平面图貌}

\begin{itemize}
    \item \textbf{等倾斜线法}：曲线 $X(x,y) = 0$（垂直等倾斜线）上 $x' = 0$，轨线方向垂直；曲线 $Y(x,y) = 0$（水平等倾斜线）上 $y' = 0$，轨线方向水平。两曲线交点为\textbf{奇点}。两条曲线之间的区域内 $x', y'$ 不变号，轨线方向在 $90^\circ$ 范围内保持不变，有 $(+,+), (+,-), (-,+), (-,-)$ 四种走向。利用等倾斜线可基本确定轨线在全相平面上的走向。
\end{itemize}

\textbf{两种群模型}：
\begin{equation}\begin{cases} \dfrac{dx}{dt} = rx(1 - ax - by) \\ \dfrac{dy}{dt} = sy(1 - cx - dy) \end{cases} \tag{6.29}\end{equation}

\begin{itemize}
    \item $b > 0, c > 0$：\textbf{竞争系统}；$b, c$ 异号：\textbf{被捕食-捕食系统}；$b < 0, c < 0$：\textbf{共生系统}
    \item 因 $x, y$ 表示种群数量，方程仅在\textbf{第一象限}有意义；坐标轴由整条轨线和奇点组成，由唯一性第一象限轨线不能穿过坐标轴
    \item \textbf{竞争系统四种结果}：\textcircled{1} $a < c, b < d$：轨线趋于 $x$ 轴上 $A$ 点，$y$ 种群绝灭；\textcircled{2} $a > c, b > d$：趋于 $y$ 轴上 $D$ 点，$x$ 种群绝灭；\textcircled{3} $a > c, b < d$：趋于 $E$ 点（结点），两种群按比例共存；\textcircled{4} $a < c, b > d$：视区域趋于 $A$、$D$ 或 $E$ 点，$E$ 点为鞍点
    \item \textbf{Volterra 被捕食-捕食模型}：相平面上可能出现极限环、焦点或中心；通解由分离变量法得首次积分式 $x^c e^{-dx}y^a e^{-by} = k$，轨线是一族围绕奇点 $E$ 的闭轨线
\end{itemize}

\textbf{分界线、同宿、异宿环（轨）}：
\begin{itemize}
    \item \textbf{分界线}：从奇点到奇点的轨线
    \item \textbf{同宿环（轨）}：一条分界线与一个奇点构成的环
    \item \textbf{异宿轨}：分界线两端是不同奇点
    \item \textbf{异宿环}：多条分界线与多个奇点构成的环
    \item 定义中奇点可换为极限环
\end{itemize}

\textbf{全局图貌}：通过\textbf{庞加莱变换}将无穷远化为原点，分析无穷远处的轨线图貌，与有限区域分析结合得到全局图貌。

\textbf{作业（6.2）}：p169-170：1（1、2）、2、3、4（1）、5（1、2）、6（1）、7（3）

\section{哈密顿系统、孤立子和混沌（概述）}

\begin{itemize}
    \item 哈密顿系统：以 Hamilton 函数为第一积分的保守系统
    \item 孤立子：非线性方程具有粒子性的稳定行波解
    \item 混沌：Lorenz 方程等表现出对初值的敏感性（蝴蝶效应，见第一章）
\end{itemize}
